\documentclass[11pt]{article}
\usepackage[T1]{fontenc}
\usepackage{textcomp}
\usepackage[margin=0.75in]{geometry}
\usepackage{amsmath,amssymb,amsthm}
\usepackage{array,booktabs}
\usepackage{hyperref}
\setlength{\parindent}{0pt}
\newtheorem{theorem}{Theorem}
\theoremstyle{definition}
\newtheorem{definition}{Definition}
\title{Flow Proof Artifact: Bool-de-morgan}
\author{Generated by \texttt{flow doc proof}}
\date{Flow Proof Book}
\begin{document}
\maketitle
De Morgan law linking negation with conjunction and disjunction.\\[0.5em]
\textbf{Source.} boole — https://en.wikipedia.org/wiki/De\_Morgan\%27s\_laws\\[0.5em]
\bigskip
\noindent{\large \textbf{Derived fact 1.} \textit{¬(a ∧ b) = (¬a) ∨ (¬b)} \hfill boolean truth values \cdot negation \cdot de Morgan for conjunction and disjunction}\par
\smallskip\noindent\textit{Coordinate: boolean truth values \cdot negation \cdot de Morgan for conjunction and disjunction}\par
\smallskip\noindent\textit{Source: boole — https://en.wikipedia.org/wiki/De\_Morgan\%27s\_laws}\par
\smallskip\noindent\textit{Built on: double negation returns the value, for negation on boolean truth values}\par
\medskip
\noindent\textbf{Goal.} ¬(a ∧ b) = (¬a) ∨ (¬b)\par
\noindent$\displaystyle \forall a \in \{\mathsf{true}, \mathsf{false}\} \forall b \in \{\mathsf{true}, \mathsf{false}\}\quad !(a \land b) = (!a) \lor (!b)$\par
\medskip
\noindent
\begin{tabular}{@{} >{\raggedright\arraybackslash}p{0.06\textwidth} >{\raggedright\arraybackslash}p{0.47\textwidth} >{\raggedleft\arraybackslash}p{0.41\textwidth} @{}}
\textbf{\#} & \textbf{Proof} & \textbf{Mathematics} \\
\midrule
\textbf{1.} & We prove that de Morgan for conjunction and disjunction for negation on boolean truth values. & \\[0.45em]
\textbf{2.} & We split into exhaustive cases — the claim must hold in each one. & \\[0.45em]
\textbf{3.} & Case 1 (see step 2): suppose a holds. & \\[0.45em]
\textbf{4.} & Case 2 (see step 2): suppose b holds. & \\[0.45em]
\textbf{5.} & From step 4, this implies !(the conjunction of a and b) equals (!a) or (!b) in this case. & $\displaystyle !(a \land b) = (!a) \lor (!b)$ \\[0.45em]
\textbf{6.} & Case 3 (see step 2): neither disjunct holds. & \\[0.45em]
\textbf{7.} & From step 6, this implies !(the conjunction of a and b) equals (!a) or (!b) in this case. & $\displaystyle !(a \land b) = (!a) \lor (!b)$ \\[0.45em]
\textbf{8.} & Case 4 (see step 2): suppose b holds. & \\[0.45em]
\textbf{9.} & From step 8, this implies !(the conjunction of a and b) equals (!a) or (!b) in this case. & $\displaystyle !(a \land b) = (!a) \lor (!b)$ \\[0.45em]
\textbf{10.} & Case 5 (see step 2): neither disjunct holds. & \\[0.45em]
\textbf{11.} & From step 10, this implies !(the conjunction of a and b) equals (!a) or (!b). Together with the other cases (step 3, step 4, step 6, step 8, and step 10), the goal is discharged. Hence proven. & $\displaystyle !(a \land b) = (!a) \lor (!b)$ \\[0.45em]
\bottomrule
\end{tabular}
\label{thm:Bool-negation-de_morgan_for_conjunction_and_disjunction}
\par\medskip\hrule
\end{document}
